\begin{solapartat}
\punts{0,25} Cal dibuixar línies de camp que surten del pol nord i van cap al pol sud.
\figura[0.55]{sol_fig1.png}

\punts{0,5} Segons la llei de Faraday, s'induirà una força electromotriu i, per tant, un corrent elèctric si el flux de camp magnètic varia.
Atès que l'imant s'acosta a la bobina, la intensitat del camp magnètic dins de la bobina augmenta i, per tant, també augmentarà el flux de camp magnètic a través de la bobina. Per tant s'induirà un corrent. Perquè la resposta sigui totalment correcta, cal esmentar la llei de Faraday i cal indicar que el flux de camp magnètic varia, no n'hi ha prou a dir que la intensitat del camp magnètic varia.

\punts{0,5} Segons la llei de Lenz, el sentit del corrent induït serà tal que el camp magnètic induït s'oposarà a la variació del flux magnètic que el genera. En aquest punt, dos raonaments són possibles:
\begin{enumerate}[label=\arabic*)]
  \item Com que la intensitat del camp magnètic dins l'espira augmenta, el camp magnètic induït tindrà un sentit oposat al camp magnètic generat per l'imant, per tant, segons la regla de la mà dreta, el corrent circularà en el sentit indicat a la figura.
  \item El corrent induït generarà un camp magnètic que crearà una força repulsiva sobre l'imant, a fi d'aturar el seu moviment. Com que el pol nord avança cap a l'espira, les línies del camp magnètic induït apuntaran cap al pol nord de l'imant, com s'indica a la figura.
\end{enumerate}
\figura[0.4]{sol_fig2.png}
\end{solapartat}

\begin{solapartat}
\punts{0,25} En aquest cas, la intensitat del camp magnètic dins l'espira disminueix i, per tant, el flux de camp magnètic decreix. Com a resultat, segons la llei de Faraday, es generarà una força electromotriu i un corrent elèctric induïts que, segons la llei de Lenz, s'oposaran al canvi.

\punts{1} Atès que ara el flux decreix, el sentit del corrent i del camp magnètic induïts serà l'oposat a l'apartat a), o també podem dir paral·lel al camp magnètic creat per l'imant, com s'indica a la figura.
\figura[0.4]{sol_fig3.png}
Donat el sentit del camp magnètic induït, la força sobre l'imant serà \textcolor{red}{atractiva}.

\destaca{Alternativament}, també es pot raonar directament a partir de la llei de Lenz; l'efecte del camp magnètic induït és aturar el moviment de l'imant, i com que ara l'imant s'allunya, la força haurà de ser \textcolor{red}{atractiva}.
\end{solapartat}
